% !TEX root = chapter-standalone.tex

\section{Product}
\label{sec:cartesian-product}

\begin{ctdefinition}[Cartesian product of two sets]
    \label{def:cartesian-product}
    \SYNDEF{cartesian product}
    Given sets~\setA and~\setB, their \emph{cartesian product} is a new set -- denoted~$\setA \cartprod \setB$ -- whose elements are precisely all possible 2-tuples~$\tup{\ela, \elb}$ such that the first entry~$\ela$ is an element of~\setA and the second entry~$\elb$ is an element of~\setB.
\end{ctdefinition}
For example, if~$\setA = \makeset{ \sapple, \sgrapes, \scarrot}$ and~$\setB = \makeset{ \stea, \swater }$, then
\begin{equation}
    \setA \cartprod \setB = \makesett{ \tup{\sapple, \stea}, \tup{\sapple, \swater}, \tup{\sgrapes, \stea}, \tup{\sgrapes, \swater}, \tup{\scarrot, \stea}, \tup{\scarrot, \swater}}.
\end{equation}

In the special case where~$\setA = \Emptyset$ or~$\setB = \Emptyset$, then~$\setA \cartprod \setB = \Emptyset$.

Another way to represent the \SY{cartesian product} is the following:
%
\equationsag{30_cartprod_repr}{eq:fig_cartprod_repr}

%The above defines the cartesian product of two sets.
%In fact we have a notion of $n$-fold cartesian product for any $n \setin \natnumbers$.
%Given sets~$\setA_1, \setA_2, \ldots, \setA_n$,
%\begin{equation}\label{eq:nfold-cart-prod}
%    \setA_1 \cartprod \setA_2 \cartprod \dots \cartprod \setA_n
%\end{equation}
%is the set of $n$-tuples~$\tup{\ela_1, \ldots, \ela_n}$ with~$\ela_i \setin\setA_i$ for all~$i = 1,2,\ldots,n$.
%
%\begin{example}
%    The vector spaces $\reals^n$ that we know and love are $n$-fold cartesian products of the set $\reals$ with itself.
%    In other words, a set of the form \ref{eq:nfold-cart-prod}, with $\setA_1, = \setA_2 = \dots = \setA_n = \reals$.
%\end{example}

\begin{remark}
    For finite sets~\setA and~\setB, the size of~$\setA \cartprod \setB$ is the product (multiplication) of the sizes of~\setA and~\setB:
    \begin{equation}
        \cardof{\setA \cartprod \setB} = \cardof\setA \cdot \cardof \setB.
    \end{equation}
    This is one reason why we think of~$\setA \cartprod \setB$ as a kind of multiplication of sets.
\end{remark}

\begin{remark}[Do you want to be more formal?]
    \label{rem:formal-cartesian}
    In formal set theory, 2-tuples (ordered pairs) are often defined by setting
    \begin{equation}
        \label{eq:ordered-pair}
        \tup{\ela, \elb} \definedas \makeset{ \makeset{\ela }, \makeset{ \ela, \elb } }.
    \end{equation}
    In this case, the \SY{cartesian product} is
    \begin{equation}
        \setA \cartprod \setB = \makeset{ \elc \setin\powerset \powerset (\setA \setunion \setB) \mid \exists \ela \setin\setA, \exists \elb \setin\setB : \elc = \tup{\ela, \elb} }.
    \end{equation}
    We will, however, treat tuples as a primitive construction (\ie, without reference to formal set theory), and hence for us also the construction of ``\SY{cartesian product of sets}'' is primitive.
\end{remark}

The definition \cref{def:cartesian-product} is for a \emph{binary} operation of cartesian product: one where there are two factors.
We can make an analogous ``n-ary'' definition for any finite number of factors.

\begin{ctdefinition}[n-ary Cartesian product of sets]
    \label{def:n-ary-cartesian-product}
    Let $n \setin \natnumbers$.
    Given sets $\setAn{1}, \setAn{2}, \dots, \setAn{n}$, their \emph{cartesian product} is the set $\setAn{1} \cartprod \setAn{2} \cartprod \dots \cartprod \setAn{n}$ whose elements are precisely all possible $n$-tuples~$\tup{\ela_{1}, \ela_{2}, \dots, \ela_{n}}$ such that each entry~$\ela_{i}$ is an element of~$\setAn{i}$, for each $i \setin \makeset{1, \dots, n}$.
\end{ctdefinition}
